\subsection{A genus-six trace and the remaining integral detection problem}

The following gives an exact same-discriminant product.  Its trace is
integral, but a unit value is a further assertion.  In particular, the
result does not provide an auxiliary field or close the arithmetic gate.

\begin{proposition}[KT-SENARY-GENUS-TRACE-DETECTOR]
Let $p\geq31$ be prime, $p\equiv3\pmod4$, and put
$b=(p+3)/2$ and $k=b+3$.  Write $E_k^{(6)}$ for the degree-six,
level-one Siegel Eisenstein series with constant coefficient $1$, and set
\begin{equation}
 F_p=\kappa_k E_k^{(6)},\qquad
 \kappa_k=\frac{\zeta(1-k)}8
                  \prod_{j=1}^3\zeta(1-2k+2j).
 \tag{SG.1}
\end{equation}
Then $\kappa_k/B_k\in\mathbb Z_{(p)}^\times$.

Let $D<0$ be odd and fundamental, put $d=|D|$, suppose $p\nmid d$,
and let $t$ be the number of distinct prime divisors of $d$.
Let $\mathcal L_D$ be the genus of positive definite even lattices of
rank six whose local lattices are
$H\oplus H\oplus\mathcal O_{\mathbb Q(\sqrt D)}(-1)$.
This genus is nonempty and its lattices have determinant $d$.
For a lattice $L$ in this genus, let $T_L$ be half its Gram matrix.
Then
\begin{equation}
 a(F_p,T_L)=-\beta_b(\chi_D),\qquad
 \mathcal T_D(F_p):=
  \sum_{[L]\in\mathcal L_D/\cong}
             \frac{a(F_p,T_L)}{|\operatorname O(L)|}
 =-\frac{\beta_b(\chi_D)\beta_3(\chi_D)}
                    {2^{t-1}\,23040}.
 \tag{SG.2}
\end{equation}
Every summand in this finite trace is $p$-integral.  Its reduction is
nonzero if and only if both Bernoulli quotients are units.  If also
$d<4p$, the already proved small-discriminant bound supplies
$p\nmid h(D)$, so a nonzero trace in that range supplies all three
units for the automatic progression under consideration.
\end{proposition}

\begin{proof}
We use two explicitly located primary inputs.  First, Kohnen's
coefficient formula is Theorem 21 on p.~22 of his author-hosted
\emph{A short course on Siegel modular forms}.  The definitions of
$D_T$, $f_T$, and $\varphi(a;T)$ are on pp.~20--21.
For degree $2n$ it states
\[
 a(E_k^{(2n)},T)=
 \frac{2^n}{\zeta(1-k)\prod_{j=1}^n\zeta(1-2k+2j)}
 \sum_{a\mid f_T}a^{k-n-1}\varphi(a;T)
 H\bigl(k-n,|D_{T,0}|(f_T/a)^2\bigr).
\]
Here $(-1)^n\det(2T)=D_{T,0}f_T^2$.  In our case $n=3$,
$D_{T,0}=D$ and $f_T=1$.  The defining sum for $\varphi(1;T)$
has the single unimodular class and gives $\varphi(1;T)=1$.
Thus (SG.1) gives
$a(F_p,T_L)=H(b,d)=L(1-b,\chi_D)=-\beta_b(\chi_D)$.
The coefficient theorem is imported; its earlier lifting proof is not
being re-proved here.

Second, the genus and its local specification are given by
Ibukiyama--Williams, \emph{Hermitian modular forms and algebraic
modular forms on $\operatorname{SO}(6)$}, Section 5, author PDF
pp.~23--26.  We give the numerical mass calculation explicitly, to
fix the normalization as the sum of reciprocals of the \emph{full}
orthogonal automorphism groups.  Their Minkowski--Siegel formula is
\[
 \operatorname{mass}(\mathcal L_D)
   =\frac{3d^{7/2}}{2\pi^9}\prod_q\alpha_q^{-1}.
\]
For odd $q\nmid d$ the local factor is
$(1-q^{-2})(1-q^{-4})(1-\chi_D(q)q^{-3})$; for $q\mid d$ it is
$2q(1-q^{-2})(1-q^{-4})$.  At $2$ the displayed local quantities
$P=(1-2^{-2})(1-2^{-4})(1-2^{-6})$ and
$E=2/(1+\chi_D(2)2^{-3})$ give
\[
 \alpha_2=2^5PE
 =2^6(1-2^{-2})(1-2^{-4})(1-\chi_D(2)2^{-3}).
\]
Consequently
\[
 \prod_q\alpha_q^{-1}
   =\frac{\zeta(2)\zeta(4)L(3,\chi_D)}{2^{t+6}d},
 \qquad
 \operatorname{mass}(\mathcal L_D)
   =\frac{3d^{5/2}\zeta(2)\zeta(4)L(3,\chi_D)}
                         {2^{t+7}\pi^9}.
\]
The primitive odd quadratic functional equation, equivalently the
Bernoulli Fourier calculation used in (QD.1), yields
$d^{5/2}L(3,\chi_D)/\pi^3=2B_{3,\chi_D}/3$.
Substituting $\zeta(2)=\pi^2/6$ and
$\zeta(4)=\pi^4/90$ therefore gives
\begin{equation}
 \operatorname{mass}(\mathcal L_D)
   =\frac{B_{3,\chi_D}}{2^{t-1}2^9 3^3 5}
   =\frac{\beta_3(\chi_D)}{2^{t-1}23040}.
 \tag{SG.3}
\end{equation}
Combining this with the constant value of the Fourier coefficient
proves (SG.2).

We check the integral assertions separately.  A finite subgroup of
$\operatorname{GL}_6(\mathbb Z)$ has order prime to $p\geq31$:
otherwise Cauchy's theorem gives an element of order $p$, whose
nontrivial rational representation has dimension at least $p-1>6$.
Thus each automorphism denominator is a $p$-unit.  Moreover
$3,b\leq p-2$ and $p\nmid d$, so the Bernoulli-polynomial formula
shows directly that both $\beta_3(\chi_D)$ and $\beta_b(\chi_D)$
are $p$-integral.  Since $23040$ and $2$ are $p$-units, the trace
criterion follows.

Finally, the three indices $2k-2j$ are $p+7,p+5,p+3$.
Kummer congruences give
\[
 \frac{B_{p+7}}{p+7}\equiv\frac{B_8}{8}=-\frac1{240},\quad
 \frac{B_{p+5}}{p+5}\equiv\frac{B_6}{6}=\frac1{252},\quad
 \frac{B_{p+3}}{p+3}\equiv\frac{B_4}{4}=-\frac1{120}
 \pmod p.
\]
All are units in the stated range, and $k<p-1$.  Hence
$\kappa_k/B_k=-(8k)^{-1}\prod_{j=1}^3\zeta(1-2k+2j)$ is a
$p$-unit.  The assertion about $h(D)$ is exactly
KT-SMALL-DISCRIMINANT-CLASS-BOUND.
\end{proof}

\paragraph{A real cancellation in this trace.}
The previously certified example
KT-JACOBI-THIRD-UNIT-NONIMPLICATION has $p=31$, $b=17$, $D=-79$,
$\beta_{17}(\chi_{-79})\equiv20\pmod{31}$ and
$\beta_3(\chi_{-79})=992$.  Formula (SG.3) gives
$\operatorname{mass}(\mathcal L_{-79})=31/720$.
Thus every individual $a(F_{31},T_L)$ is a unit, while the integral
genus trace is zero modulo $31$.  Nonvanishing of an individual
full-rank Siegel coefficient therefore does not force this trace to
be a unit, even for the actual Eisenstein series used here.

\begin{lemma}[KT-ANALYTIC-RESIDUE-NONSATURATION]
An analytically computed unit residue of a Dirichlet series need not
detect a unit among its $p$-integral coefficients, even when every
index is prime to $p$.
\end{lemma}
\begin{proof}
The absolutely convergent series on $\Re s>1$
\[
 A_p(s)=p\sum_{p\nmid n}n^{-s}=p(1-p^{-s})\zeta(s)
\]
has every coefficient in $p\mathbb Z$, but its meromorphic residue
at $s=1$ is $p-1$, a $p$-unit.  A residue argument for the trace above
must therefore establish an integral coefficient functional in that
specific modular setting; complex continuation and the residue value
alone do not establish it.
\end{proof}

\paragraph{Primary locators and normalization caution.}
Kohnen's primary author notes are at
\url{https://www.mathi.uni-heidelberg.de/~winfried/siegel2.pdf},
pp.~20--22, especially Theorem 21.  Ibukiyama--Williams is published
in J. Number Theory \textbf{286} (2026), 230--279,
DOI 10.1016/j.jnt.2026.02.009; the checked author PDF is
\url{https://btw-47.github.io/hermitian_orthogonal.pdf}.
Its Proposition 12 statement on p.~24 prints $3^2$ in the mass
denominator, while the proof conclusions on pp.~25--26 have $3^3$.
The explicit calculation above confirms $3^3$; for $D=-3$ it gives
the $E_6$ mass $1/103680$.  We do not use the numerical prefactor of
the 1996 Ibukiyama--Katsurada Koecher--Maa\ss{} announcement, or
assert that an unexamined analytic residue supplies an integral unit
functional.  No determinant bound or unit trace is proved here.
